\section{Modelo de datos}
\label{sec:modelo}

El esquema \texttt{public} agrupa $17$ tablas crudas en cinco bloques:
\emph{producto} (\texttt{product},\allowbreak{} \texttt{central\_product},\allowbreak{}
\texttt{category}, \texttt{brand}, \texttt{offer},\allowbreak{}
\texttt{product\_offer}); \emph{almacén} (\texttt{warehouse},\allowbreak{}
\texttt{warehouse\_location},\allowbreak{} \texttt{central\_inventory},\allowbreak{}
\texttt{inventory}); \emph{tienda y geografía} (\texttt{store},
\texttt{city\_zone}); \emph{cliente y ventas} (\texttt{customer},
\texttt{sale}, \texttt{sale\_item}); y \emph{devoluciones}
(\texttt{return\_item}, \texttt{return\_reason}). Las relaciones clave
son \texttt{sale}\,1--N\,\texttt{sale\_item}, \texttt{sale}\,N--1
\texttt{customer} y \texttt{store}, \texttt{sale\_item}\,N--1
\texttt{product}, y \texttt{sale\_item}\,1--N\,\texttt{return\_item}
(figura~\ref{fig:er}).

\begin{figure}[H]
  \centering
  \includegraphics[width=0.92\linewidth]{er_diagram.png}
  \caption{Esquema \texttt{public}: $17$ tablas agrupadas en cinco bloques
  funcionales.}
  \label{fig:er}
\end{figure}

El esquema \texttt{dwh} sigue un patrón clásico en estrella con seis
dimensiones (\texttt{dim\_date}, \texttt{dim\_customer},
\texttt{dim\_product}, \texttt{dim\_store}, \texttt{dim\_offer},
\texttt{dim\_return\_reason}) y dos hechos: \texttt{fact\_sales}
(grano $=$ línea de venta, $42\,555$~filas, métricas \texttt{quantity},
\texttt{unit\_price}, \texttt{unit\_cost}, \texttt{subtotal},
\texttt{cost}, \texttt{margin}) y \texttt{fact\_returns}
(grano $=$ línea devuelta, $2\,330$~filas, métricas \texttt{quantity} y
\texttt{returned\_value}). Véase la figura~\ref{fig:dim}.

\begin{figure}[H]
  \centering
  \includegraphics[width=0.92\linewidth]{dim_diagram.png}
  \caption{Esquema \texttt{dwh}: seis dimensiones más dos hechos. La línea
  discontinua representa el \emph{drill-through} de
  \texttt{fact\_returns} hacia \texttt{fact\_sales} por
  \texttt{sale\_item\_id}.}
  \label{fig:dim}
\end{figure}

\subsection*{Decisiones de diseño}

Todas las dimensiones tienen una \emph{surrogate key} \texttt{*\_sk}
(\texttt{serial}) y conservan la clave de negocio como \texttt{unique},
lo que hace la carga idempotente (\texttt{truncate \dots restart
identity}). \texttt{dim\_product} se construye con \texttt{LEFT JOIN}
entre \texttt{product} ($50$~filas) y \texttt{central\_product}
($49$~filas) — un producto carece de coste y su margen se gestiona con
\texttt{coalesce(unit\_cost,\,0)}. \texttt{fact\_returns} desnormaliza
\texttt{customer\_sk}, \texttt{product\_sk} y \texttt{store\_sk} desde
\texttt{fact\_sales} a través de \texttt{sale\_item\_id}, lo que evita
\emph{re-joins} en analítica de devoluciones. En el área \texttt{staging}
se recomputa \texttt{subtotal} en el $100\%$ de las líneas porque $8$
filas raw eran inconsistentes con $\texttt{quantity}\times
\texttt{unit\_price}$; el test
\texttt{test\_revenue\_matches\_between\_staging\_and\_dwh} certifica
que la corrección no rompe la conservación de ingresos. El DDL completo
está en \texttt{sql/02\_dwh\_schema/}.
